% =============================================================================
% DOCUMENT: Thesis Template - Chapter 1: Customer Lifetime Value Foundations
% Institution: ENSSEA - École Nationale Supérieure de Statistique et d'Économie Appliquée
% Format: Master's Memoir - Following ENSSEA Guidelines
% =============================================================================

\documentclass[12pt,a4paper,oneside]{report}

% ============================================================================
% PACKAGES AND CONFIGURATION
% ============================================================================

\usepackage[utf-8]{inputenc}
\usepackage[T1]{fontenc}
\usepackage[english]{babel}
\usepackage{geometry}
\usepackage{setspace}
\usepackage{amsmath}
\usepackage{amssymb}
\usepackage{graphicx}
\usepackage{hyperref}
\usepackage{booktabs}
\usepackage{array}
\usepackage{float}
\usepackage{fancyhdr}
\usepackage{titlesec}
\usepackage{xcolor}

% ============================================================================
% PAGE GEOMETRY (ENSSEA Standards)
% ============================================================================

\geometry{
    left=3cm,
    right=2cm,
    top=2cm,
    bottom=2cm,
    headheight=1cm,
    headsep=0.5cm
}

% ============================================================================
% TYPOGRAPHY AND FORMATTING (ENSSEA Standards)
% ============================================================================

% Chapter titles formatting
\titleformat{\chapter}[hang]
{\normalfont\bfseries\fontsize{20}{24}\selectfont}
{\thechapter.}{1em}{}

% Section formatting (first-level headings)
\titleformat{\section}[hang]
{\normalfont\bfseries\fontsize{14}{16}\selectfont}
{\thesection}{1em}{}

% Subsection formatting (second-level headings)
\titleformat{\subsection}[hang]
{\normalfont\bfseries\fontsize{12}{14}\selectfont}
{\thesubsection}{1em}{}

% Subsubsection formatting (third-level headings)
\titleformat{\subsubsection}[hang]
{\normalfont\bfseries\fontsize{12}{14}\selectfont\underline}
{\thesubsubsection}{1em}{}

% Line spacing: 1.5
\onehalfspacing

% Paragraph formatting: justified with no indent, paragraph spacing
\setlength{\parindent}{0pt}
\setlength{\parskip}{0.6cm}

% ============================================================================
% HEADER AND FOOTER (ENSSEA Standards)
% ============================================================================

\pagestyle{fancy}
\fancyhf{}
\fancyfoot[C]{\thepage}
\renewcommand{\headrulewidth}{0pt}
\renewcommand{\footrulewidth}{0pt}

% ============================================================================
% DOCUMENT BEGINS
% ============================================================================

\begin{document}

% Suppress page numbers for front matter
\pagenumbering{roman}

% ============================================================================
% CHAPTER 1: THEORETICAL FRAMEWORKS OF CUSTOMER LIFETIME VALUE
% ============================================================================

\chapter{Theoretical Foundations of Customer Value: Conceptual Framework and Probabilistic Models}

% Introduction section of Chapter 1
\section{1.1 - General Introduction to Customer Value Concepts}

Effective customer relationship management (CRM) constitutes a major strategic challenge for contemporary organizations. At the heart of this issue lies a fundamental question: how can one rigorously and predictively evaluate the value that a customer represents for the enterprise? This question has prompted practitioners and researchers to develop sophisticated conceptual and methodological frameworks for measuring what is commonly termed Customer Lifetime Value (CLV).

Traditionally, customer segmentation approaches relied on relatively simple demographic or geographic criteria. However, the evolution of computational capabilities and the increasing availability of detailed transactional data have enabled the development of more granular and behavioral approaches. These developments have led to progressive understanding that certain customers, while fewer in number, can generate disproportionate value compared to other segments of the customer base.

The concept of CLV has thus emerged as a central pillar of customer relationship management strategy, enabling companies to identify, quantify, and act upon the heterogeneity that characterizes their customer base. More specifically, CLV allows practitioners to address the following strategic questions: Which customers are worth the most? How should resources be allocated across the customer base to maximize long-term profitability? What is the economic justification for retention investments?

% ============================================================================
% SECTION 1.2
% ============================================================================

\section{1.2 - Theoretical Frameworks and Probabilistic Models of Customer Lifetime Value}

\subsection{1.2.1 - Conceptual Definition and Economic Foundations of CLV}

Customer Lifetime Value represents the present value of expected future cash flows attributable to the customer-company relationship throughout its duration. Formally, CLV can be expressed as:

\begin{equation}
CLV = \sum_{t=0}^{T} \frac{(R_t - C_t)}{(1+d)^t}
\label{eq:clv_general}
\end{equation}

where:
\begin{itemize}
\item $R_t$ represents revenues generated by the customer in period $t$,
\item $C_t$ represents costs incurred to serve the customer in period $t$,
\item $d$ represents the discount rate reflecting the time value of money,
\item $T$ represents the expected duration of the customer relationship.
\end{itemize}

This formulation, grounded in microeconomic theory and the principles of net present value, provides a formal framework within which CLV can be estimated and operationalized. The key challenge in implementing this framework resides in the estimation of future cash flows $(R_t - C_t)$ and the appropriate selection of the temporal horizon $T$.

CLV serves as both a descriptive metric and a decision-making tool. As a descriptive metric, it provides an assessment of the current economic value attributable to a customer or customer segment. As a decision-making tool, it enables enterprises to allocate resources more efficiently by concentrating acquisition and retention efforts on customer segments demonstrating the highest CLV.

\subsection{1.2.2 - Probabilistic and Econometric Models for CLV Prediction}

The modeling of CLV has led to the development of several classes of models, each presenting specific advantages and limitations. These approaches can be classified according to their theoretical foundation: models based on Recency, Frequency, and Monetary value (RFM) analysis, contractual probabilistic models, and non-contractual probabilistic models.

\subsubsection{1.2.2.1 - RFM Models and Descriptive Approaches}

The RFM (Recency, Frequency, Monetary) approach represents a heuristic method widely used in practice for segmenting and evaluating customers\footnote{Hughes, A.M. (1996). "The Customer Loyalty Solution: What Works (and What Doesn't) in Customer Loyalty Programs". New York: McGraw-Hill.}. This method is based on the hypothesis that three variables—recency (R) of the last purchase, frequency (F) of prior purchases, and amount (M) spent—constitute good predictors of future purchase behavior and thus the customer's potential value.

Although simple to implement, the RFM approach presents significant limitations. Notably, it does not provide rigorous probabilistic estimates of CLV and implicitly assumes a linear relationship between these three variables and future value, which may not be empirically verified.

\subsubsection{1.2.2.2 - Probabilistic Models in Contractual Contexts}

In contexts where customer relationships present an explicit contractual structure (subscriptions, service contracts), probabilistic models based on the theory of Poisson processes and Markov chains have proven particularly effective\footnote{Soufiani, H., Jedidi, K., \& Degeratu, A.M. (2012). "A model of simultaneous diagonal and off-diagonal asymmetries in consumer preferences". Marketing Science, 31(6), 918-940.}.

These models formalize the process of customer churn as a stochastic process governed by a constant or time-varying intensity parameter. A fundamental example is the Exponential-Gamma (EG) model, which assumes that churn rates are distributed according to a Gamma distribution among customers, while the moments of churn for a given customer follow a Poisson process.

\subsubsection{1.2.2.3 - Non-contractual Probabilistic Models: The BG/NBD}

In non-contractual contexts, where customers are not formally bound by a contract and where churn is not directly observable, the estimation of CLV becomes substantially more complex. The BG/NBD (Beta-Geometric/Negative Binomial Distribution) model, developed by Fader, Hardie, and Lee, revolutionized this problem by providing a rigorous and parsimonious probabilistic framework\footnote{Fader, P.S., Hardie, B.G.S., \& Lee, K.L. (2005). "Counting Your Customers the Easy Way: An Alternative to the Pareto/NBD Model". Journal of Marketing Research, 42(2), 139-155.}.

The BG/NBD model rests on several fundamental assumptions\footnote{Fader, P.S., \& Hardie, B.G.S. (2009). "Probability Models for Customer-Base Analysis". Journal of Interactive Marketing, 23(1), 61-69.} :

\begin{enumerate}
\item As long as a customer is active, the number of transactions during a specified period follows a Poisson distribution with rate $\lambda$ (heterogeneous among customers).

\item Each customer possesses an unconditional probability $p$ of becoming inactive after each transaction, independent of the timing and frequency of prior transactions.

\item Heterogeneity within the population regarding the transaction rate $\lambda$ follows a Gamma distribution.

\item Heterogeneity within the population regarding the churn probability $p$ follows a Beta-Geometric distribution.

\item The parameters $\lambda$ and $p$ are independently distributed among customers.
\end{enumerate}

The BG/NBD model offers several notable advantages. First, it does not require cross-sectional observational data, making it applicable to many practical contexts. Second, its parameters admit direct behavioral interpretation. Third, it has demonstrated excellent predictive performance in practical applications\footnote{Fader, P.S., Hardie, B.G.S., \& Lee, K.L. (2003). "Counting Your Customers the Easy Way: An Alternative to the Pareto/NBD Model". Working paper, The Wharton School.}.

\subsubsection{1.2.2.4 - Gamma-Gamma Models for Monetary Value}

Parallel to the development of models for predicting transaction numbers, particular attention has been devoted to modeling the monetary value of transactions. The Gamma-Gamma model constitutes the standard approach for this dimension\footnote{Colombo, R., \& Morrison, D.E. (1989). "A Brand Switching Model with Implications for Marketing Expenditures". Marketing Science, 8(1), 89-99.}.

This model assumes that the average monetary value of each transaction for a given customer is represented by a Gamma distribution, with a heterogeneity structure also Gamma at the population level. Formally, the joint distribution of the expected monetary value $m_x$ and the number of observed transactions $x$ is given by:

\begin{equation}
E[M|x, m_x] = \frac{(r+1)m_x}{r}
\label{eq:gamma_gamma}
\end{equation}

where $r$ represents the shape parameter of the Gamma distribution at the population level.

\subsection{1.2.3 - Relationship Between CLV, Customer Retention, and Profitability}

The establishment of rigorous relationships between CLV, retention strategies, and realized profitability constitutes a central issue for strategic customer relationship management. Several empirical studies have documented that the relationship between estimated CLV and realized profitability is not always straightforward and depends crucially on cost structure\footnote{Reinartz, W.J., Thomas, J.S., \& Kumar, V. (2005). "Balancing Acquisition and Retention Resources to Maximize Customer Profitability". Journal of Marketing, 69(1), 63-79.}.

Customer retention, traditionally conceived as an objective in itself, turns out to be justified only if it generates sufficient additional value to offset the additional costs incurred. Indeed, the relationship between the retention rate and profitability proves to be a non-monotonic function: an increase in the retention rate is not systematically associated with increased profitability, particularly when retention costs are high\footnote{Reinartz, W.J., \& Kumar, V. (2000). "On the profitability of long-life customers in a noncontractual setting: An empirical investigation and implications for marketing". Journal of Marketing, 64(4), 17-35.}.

Consequently, rational use of CLV to guide resource allocation decisions requires a nuanced understanding of this relationship. Firms must identify customer segments for which an increase in retention effectively generates an increase in CLV, and concentrate retention efforts on these segments\footnote{Pfeifer, P.E., Haskins, M.E., \& Conroy, R.M. (2005). "Customer Lifetime Value, Customer Profitability, and the Treatment of Acquisition Spending". Journal of Managerial Issues, 17(1), 11-25.}.

% ============================================================================
% CONCLUSION OF CHAPTER 1
% ============================================================================

\section{Conclusion of Chapter 1}

This chapter has presented the theoretical and conceptual foundations of customer lifetime value (CLV). We have established that CLV represents far more than a simple business metric; it constitutes an analytical framework enabling alignment of customer management decisions with long-term value creation for the enterprise.

The probabilistic and econometric models presented offer sophisticated tools for estimating CLV in varied empirical contexts, each presenting a set of advantages and limitations that must be considered when selecting a methodological approach. The BG/NBD model, in particular, has emerged as a de facto standard for non-contractual contexts due to its parsimony, rigorous theoretical foundation, and excellent predictive performance.

Finally, we have emphasized that effective use of CLV to guide strategic decisions requires nuanced understanding of the relationships between CLV, customer retention, and profitability. This understanding will be deepened in subsequent chapters, where we will examine empirical applications of these concepts and recent methodological developments in the field.

% ============================================================================
% BIBLIOGRAPHY / REFERENCES
% ============================================================================

\newpage
\section*{Chapter 1 Bibliographic References}

\begin{thebibliography}{99}

\bibitem[Colombo \& Morrison (1989)]{colombo1989} Colombo, R., \& Morrison, D.E. (1989). "A Brand Switching Model with Implications for Marketing Expenditures". \textit{Marketing Science}, 8(1), 89-99.

\bibitem[Fader et al. (2003)]{fader2003} Fader, P.S., Hardie, B.G.S., \& Lee, K.L. (2003). "Counting Your Customers the Easy Way: An Alternative to the Pareto/NBD Model". Working paper, The Wharton School.

\bibitem[Fader et al. (2005)]{fader2005} Fader, P.S., Hardie, B.G.S., \& Lee, K.L. (2005). "Counting Your Customers the Easy Way: An Alternative to the Pareto/NBD Model". \textit{Journal of Marketing Research}, 42(2), 139-155.

\bibitem[Fader \& Hardie (2009)]{fader2009} Fader, P.S., \& Hardie, B.G.S. (2009). "Probability Models for Customer-Base Analysis". \textit{Journal of Interactive Marketing}, 23(1), 61-69.

\bibitem[Gupta et al. (2004)]{gupta2004} Gupta, S., Lehmann, D.R., \& Stuart, J.A. (2004). "Valuing customers". \textit{Journal of Marketing Research}, 41(1), 7-18.

\bibitem[Hughes (1996)]{hughes1996} Hughes, A.M. (1996). \textit{The Customer Loyalty Solution: What Works (and What Doesn't) in Customer Loyalty Programs}. New York: McGraw-Hill.

\bibitem[Pfeifer et al. (2005)]{pfeifer2005} Pfeifer, P.E., Haskins, M.E., \& Conroy, R.M. (2005). "Customer Lifetime Value, Customer Profitability, and the Treatment of Acquisition Spending". \textit{Journal of Managerial Issues}, 17(1), 11-25.

\bibitem[Pfeifer \& Carraway (2000)]{pfeifer2000} Pfeifer, P.E., \& Carraway, R.L. (2000). "Modeling customer relationships as Markov chains". \textit{Journal of Interactive Marketing}, 14(2), 43-55.

\bibitem[Reinartz \& Kumar (2000)]{reinartz2000} Reinartz, W.J., \& Kumar, V. (2000). "On the profitability of long-life customers in a noncontractual setting: An empirical investigation and implications for marketing". \textit{Journal of Marketing}, 64(4), 17-35.

\bibitem[Reinartz et al. (2005)]{reinartz2005} Reinartz, W.J., Thomas, J.S., \& Kumar, V. (2005). "Balancing Acquisition and Retention Resources to Maximize Customer Profitability". \textit{Journal of Marketing}, 69(1), 63-79.

\bibitem[Soufiani et al. (2012)]{soufiani2012} Soufiani, H., Jedidi, K., \& Degeratu, A.M. (2012). "A model of simultaneous diagonal and off-diagonal asymmetries in consumer preferences". \textit{Marketing Science}, 31(6), 918-940.

\end{thebibliography}

\end{document}
